%% ma: L12-B19-0007
%% lop: 12
%% chuong: 6
%% bai: 19
%% dang: 12.19.1
%% loai: DS
%% muc: [NB, NB, TH, VD]
%% dap_an: "SĐSS"
%% nguon: "(NBV)-ĐỀ SỐ 3-MỖI NGÀY 1 ĐỀ THI 2025.pdf (D03, MỖI NGÀY 1 ĐỀ - Đề số 3), Phần II Câu 2"
%% kien_thuc: "Bài 18 (xác suất có điều kiện), Bài 19 (công thức xác suất toàn phần)"
%% trang_thai: chinh-thuc
%% ghi_chu: "Đáp án gốc S Đ S S đúng (P(B)=0,2775). Bỏ ảnh minh họa. Gần ý tưởng với 12.19.1; tách riêng vì bối cảnh y tế, có thể gộp khi duyệt. Giáo viên duyệt gộp dạng 12.19.4 vào 12.19.1 (10/10/2026)"
\begin{ex}
Khi điều tra sức khoẻ nhiều trẻ em trong lứa tuổi mầm non ở một địa phương, người ta thấy rằng có $5\%$ trẻ em bị chứng tự kỉ. Bên cạnh đó, số trẻ em thừa cân béo phì trong những trẻ em tự kỉ chiếm tới $80\%$, còn số trẻ em thừa cân béo phì trong những trẻ em không tự kỉ chỉ chiếm $25\%$. Chọn ngẫu nhiên một em bé ở địa phương trên.
\choiceTF
{Xác suất chọn được em bé tự kỉ là $0{,}5$.}
{\True Xác suất chọn được em bé thừa cân béo phì, biết em đó đã mắc chứng tự kỉ là $0{,}8$.}
{Xác suất chọn được em bé thừa cân béo phì, biết em đó không mắc chứng tự kỉ là $0{,}75$.}
{Xác suất chọn được em bé thừa cân béo phì là $0{,}2$.}
\loigiai{Gọi $A$ là biến cố ``em bé tự kỉ'', $B$ là biến cố ``em bé thừa cân béo phì''. Ta có $P(A)=0{,}05$, $P(B\mid A)=0{,}8$, $P(B\mid\overline A)=0{,}25$.
a) $P(A)=0{,}05\ne0{,}5$. Sai.
b) $P(B\mid A)=0{,}8$. Đúng.
c) $P(B\mid\overline A)=0{,}25\ne0{,}75$. Sai.
d) $P(B)=P(A)P(B\mid A)+P(\overline A)P(B\mid\overline A)=0{,}05\cdot0{,}8+0{,}95\cdot0{,}25=0{,}2775\ne0{,}2$. Sai.}
\end{ex}
